\documentclass[11pt,a4paper]{article}
\usepackage{gtu-el-yazisi}
\begin{document}
\notbilgi{Görüntü İşleme}{Hough dönüşümü ve bağlı bileşenler}{Ders9.pdf}

\sayfa{1}
\section{Hough Transform}
Input : binary image $(0,1)$ \quad [edge map]

\elcizimi[0.35]{sekiller/Ders9-s1-1.png}[1]

$y_i=a\,x_i+b$ \hfill $b$ 0 olursa $y$ merkezden geçer

\[ b=-a\,x_i+y_i \]

\elcizimi[0.4]{sekiller/Ders9-s1-2.png}[1]

doğru $a$'yı bildiğinde aynı $b$'de yoğunlaşma olacak

$x_i=1,\ y_i=1$
\begin{align*}
b&=-a\,x_i+y_i\\
b&=-a+1\\
a=-2&\Rightarrow b=3\\
a=-1&\Rightarrow b=2
\end{align*}
\begin{center}
\begin{tabular}{r|r}
$a$ & $b$\\\hline
$-2$ & $3$\\
$-1$ & $2$\\
$0$ & $1$\\
$1$ & $0$\\
$2$ & $-1$\\
$3$ & $-2$
\end{tabular}
\end{center}

\section{Örnek:}
\elcizimi[0.5]{sekiller/Ders9-s1-3.png}[1]

2 tane line var

\elcizimi[0.5]{sekiller/Ders9-s1-4.png}[1]

{\color{murekkep}
$x_i=2,\ y_i=3$
\begin{align*}
b&=-a\,x_i+y_i\\
b&=-2a+3
\end{align*}}

{\color{yesilkalem}
$x_i=4,\ y_i=1$
\[ b=-4a+1 \]}

$x_i=5,\ y_i=0$
\[ b=-5a \]

{\color{turuncukalem}
$x_i=2,\ y_i=2$
\begin{align*}
b&=-a\,x_i+y_i\\
b&=-2a+2
\end{align*}
\begin{tabular}{r|r}
$a$ & $b$\\\hline
$-2$ & $6$\\
$-1$ & $4$\\
$0$ & $2$\\
$1$ & $0$\\
$2$ & $-2$\\
$3$ & $-4$
\end{tabular}}

{\color{morkalem}
$x_i=4,\ y_i=4$
\begin{align*}
b&=-a\,x_i+y_i\\
b&=-4a+4
\end{align*}
\begin{tabular}{r|r}
$a$ & $b$\\\hline
$-2$ & $12$\\
$-1$ & $8$\\
$0$ & $4$\\
$1$ & $0$\\
$2$ & $-4$
\end{tabular}}

\sayfa{2}
eğimi 3 yaparsak 3'ün üstü line olur

\elcizimi[0.4]{sekiller/Ders9-s2-1.png}[2]

$a\to-\infty,\infty$ sıkıntı

peki hangi sıklıkla

$a$ lineer olsa da ona göre açı lineer değil

\elcizimi[0.5]{sekiller/Ders9-s2-2.png}[2]

{\color{turuncukalem}$\rho=x\cos\theta+y\sin\theta$}

$y=ax+b$

\elcizimi[0.35]{sekiller/Ders9-s2-3.png}[2]

$\rho$ ve $\theta$ aşağıdaki gibi

$\theta$ $-90\div 90$ aralığında

Fig 10.32

\sayfa{3}
\elcizimi[0.5]{sekiller/Ders9-s3-1.png}[3]

\sekilnotu[3]{Bu sayfa basılı materyal içeriyor (Şekil 10.33 ve 10.34: Hough dönüşümü parametre uzayı örnekleri); dijitalleştirilmedi.}

{\color{kirmizikalem}$\rho=y\sin\theta$}

tüm $\theta$'larda $\rho=0$

$-90$ \qquad $90$

$\leftarrow 90^\circ$

\begin{dikkat}
Sınavda çıkar
\end{dikkat}

\sayfa{4}
Daire için
\[ (x-c_1)^2+(y-c_2)^2=c_3^2 \]

\section{Connected Component}
\elcizimi[0.4]{sekiller/Ders9-s4-1.png}[4]

Haralick

\elcizimi[0.4]{sekiller/Ders9-s4-2.png}[4]

\elcizimi[0.45]{sekiller/Ders9-s4-3.png}[4]

\elcizimi[0.8]{sekiller/Ders9-s4-4.png}[4]

\elcizimi[0.4]{sekiller/Ders9-s4-5.png}[4]

\sayfa{5}
\elcizimi[0.4]{sekiller/Ders9-s5-1.png}[5]

\section{Recursive Connected Component labeling}
\noindent while end of image is reached\\
\hspace*{1.5em}do\{ continue to scan image until\\
\hspace*{4em}an unlabeled ``1'' pixel is reached.\\
\hspace*{3em}if (not end of image)\{\\
\hspace*{6em}assign a new label to the pixel\\
\hspace*{6em}recursively assign same label to\\
\hspace*{6em}``1'' neighbors \}

\elcizimi[0.35]{sekiller/Ders9-s5-2.png}[5]

soldan sağa\\
yukarıdan aşağıya

\elcizimi[0.45]{sekiller/Ders9-s5-3.png}[5]

\elcizimi[0.4]{sekiller/Ders9-s5-4.png}[5]

\end{document}
